\section{Question, contract and findings}
\label{observation-selection}

This study asks whether
\emph{acquiring} one further legal observation helps resolve an
uncertainty: a future target outcome left ambiguous by the available
history. It is a separate experiment with deterministic actors, a
covering class of deterministic systems with at most three states and
binary inputs and outputs, and two permitted queries. A conditional
target certificate is issued only when every candidate consistent with
the history and the observed query response predicts the same target.
Class coverage is an explicit premise of this certificate. From the
eight-input history's cutoff, $q_0$ applies input 0 and $q_1$ applies
input 1; each returns only that transition's output. The target is the
pair of second outputs on input sequences $(0,0)$ and $(1,0)$, each
starting at the original cutoff. Appendix~\ref{app:observation-selection}
defines the state and output conventions.

The fixed policy always chooses the first query. The development-selected
policy maximizes the number of nonempty response branches having a
single target outcome, breaks ties by minimizing the largest feasible
branch's number of distinct targets, and finally chooses the first
query. Branch and candidate counts are not calibrated probabilities.
Both arms receive exactly one observation on every draw, including
initially unambiguous and duplicate records. Separate processes and
state copies prevent cross-arm reuse. Neither actor receives the true
system's identity, private state or target answer. All predictions were
committed before a separate privileged target audit.

The protocol fixed the policy, tie rules, $N=5{,}000$, confidence
procedure and validity gate before the main draw. These are local, hash-bound
precommitments, not an independently timestamped public preregistration.
The hypothesis class covers at most three states; the sampled worlds instead
have exactly three minimal reachable states, after development exclusions.
Sampling was uniform
with replacement over 8,372 eligible behavior classes crossed with
256 input histories: 2,143,232 class--history pairs. Classes used in the
specified development samples were excluded; visible-history overlap
was permitted. The realized sample contains 4,993 distinct pairs and
3,791 distinct classes. Draws, rather than unique worlds or live model
attempts, are the inferential units. Appendix~\ref{app:observation-selection}
records the development history, sampling and statistical procedure.

\begin{table}[t]
\centering
\small
\begin{tabular}{lrrr}
\toprule
 & Fixed: certificate & Fixed: no certificate & Total\\
\midrule
Selected: certificate & 282 & 90 & 372\\
Selected: no certificate & 1 & 4,627 & 4,628\\
\midrule
Total & 283 & 4,717 & 5,000\\
\bottomrule
\end{tabular}
\caption{Separate deterministic acquisition study: certificate issuance
on 5,000 paired draws, not P/L/U feasibility labels. All 655 issued
certificates matched their targets; the 282 jointly certified draws
contribute two certificates each. That cell includes 97 draws whose
targets were already unique before the observation. The remaining
4,903 draws yield 186 new certificates in the fixed arm and 275 in
the selected arm. The net gain is $90-1=89$ on the primary $N=5{,}000$.
No certificate means unresolved under the corresponding policy, not an
impossibility proof for every one-observation method.}
\label{tab:observation-selection}
\end{table}


All 5,000 pairs completed, with no technical missingness or unknown work
counters. The selected policy issued 372 certificates (7.44\%), versus
283 (5.66\%) for the fixed policy. There were 90 selected-only and one
fixed-only certificates: 89 net additional certificates, or
\textbf{1.78 percentage points}. The pre-specified, at-least-95\%
confidence interval is \textbf{[1.23, 2.29] percentage points}.
It excludes zero and supports the narrow positive-difference claim
in this frame. It includes the two-percentage-point research-interest
threshold, so a gain of at least that size remains unresolved. The
planning power calculation is not a guarantee or a post hoc estimate
of achieved power. Of the 5,000 draws, 97 already had a unique target
before the extra observation and remain in both arms' totals. Among the
4,903 initially ambiguous draws, the observation newly certified 186
targets in the fixed arm and 275 in the selected arm; their net difference
is again 89. The primary denominator remains all 5,000 draws.

All 655 issued certificates matched the target audit; any contradiction
would have invalidated the certificate result regardless of its
interval. Nevertheless, 4,628 selected-policy draws (92.56\%) remain
uncertified. The one fixed-only certificate also rules out per-record
dominance. Query choices differ on 686 draws. Of 458 draws with a
pre-observation possible gain, only 90 realize it; the possible-loss
draw realizes one loss. These counts separate a favorable candidate
branch from the actual response. They are descriptive diagnostics,
not further primary tests or calibrated response probabilities.

Equal observation budgets did not imply equal computation. Each arm
made 5,000 environment queries and 233,280,000 candidate--history
checks. Relative to the fixed arm, the selected policy made 5,283,246
additional query--candidate evaluations and 20,000 selection-branch
checks: approximately 59,362 and 225, respectively, per \emph{net}
additional certificate. These work categories are reported separately,
not added as if they were disjoint CPU instructions. Within-arm reuse
is recorded separately and is not counted as additional work.
Actor CPU time was 136.38 versus 137.96 seconds; the fixed arm always
ran first, so timing is descriptive for this implementation and host.
No economic-utility or deployment conclusion is drawn.

This realizes the acquisition benefit bounded by the earlier branch scan.
The earlier failed prediction criteria (Appendix A.1) and the separate
decision-record crossing results remain unchanged.
No general crossing benefit, improved LLM experimentation, universal coverage
or cross-contract transfer follows.

\paragraph{Post hoc query comparators on the same inputs.}
A later reader analysis reconstructed both allowed responses on the original
5,000 draws. Always querying $q_1$ yields 219 certificates, versus 283 for
the pre-specified $q_0$ and 372 for the selected policy. A policy choosing
each query with probability $1/2$ independently of the record has an expected
$(283+219)/2=251$ certificates on these inputs. This is an exact expectation,
not an observed randomized run. Reconstructed $q_0$/$q_1$ certificates have
zero target disagreements. These comparisons are retrospective and use no
new sample or live observation; they neither replace the primary 89-net
gain and its interval nor establish optimality among other selectors.
The reader replay computes them from the reconstructed histories, candidate
supports and original world tables, separately from historical work counters.

\section{Position and scope}
Informative query selection is established \citep{cohn1994active}, including
identifiability-directed acquisition \citep{shetty2026idea}. In active automata
learning, Angluin's membership and equivalence-query setting aims to identify
a regular language \citep{angluin1987learning}; here the known covering class
is a finite transducer class, there is no equivalence oracle, and one of two
one-step queries serves a particular future target rather than full system
identification. Query by committee uses disagreement to select informative
queries \citep{seung1992query}. Our frozen branch-purity selector neither
samples a committee nor interprets candidate multiplicities as probabilities;
we make no learning-rate comparison to that method. Our contribution
is this frozen policy comparison, with realized branches, incorrect-certificate
checks and compute accounting. Neither query selection nor unanimity-based
inference in a covering finite class is claimed as novel.

The positive average difference is limited to the pre-specified finite frame;
new classes, query contracts and live LLM selection remain unassessed. The
primary comparison establishes superiority to the frozen fixed $q_0$
baseline. The $q_1$ and fair-random expectation above are separate post hoc
descriptions, not pre-specified confirmatory tests of a wider optimality claim.
The reader replay reconstructs the original sample, histories, query responses,
candidate computations and target checks from disclosed synthetic inputs,
then recovers the counts and primary interval. It collects no new sample
and does not independently attest the historical run. This result and the crossing
case study retain separate denominators and claims.
